% \input{\pSections "sec-seismology"}

\section{Seismology, Quickly}
%
%%   %%   %%   %%   %%   %%   %%   %%   %%
\subsection{Underground Explosions\jumpLittle}
\begin{frame}\frametitle{Underground Explosion B}
\center
	\href{https://apps.dtic.mil/sti/pdfs/ADA627744.pdf\#page=41}{
	\begin{overpic}[ scale = 0.5 ]
		{\pLocalGraphics ug-shot}
		\put(65, -5)	{\href{https://apps.dtic.mil/sti/pdfs/ADA627744.pdf}{\scriptsize{\colorbox{gray!10}{\cite{adushkin2007underground}}}}}
	\end{overpic}}
\end{frame}
%
%%   %%   %%   %%   %%   %%   %%   %%   %%
\subsection{Seismic Waves}
\begin{frame}\frametitle{\href{https://en.wikipedia.org/wiki/Seismic_wave}{Seismic Wave Types}}
\center
	\href{https://en.wikipedia.org/wiki/Seismic_wave}{
	\begin{overpic}[ scale = 0.25 ]
		{\pLocalGraphics seismograph/Overview_Seismic_Waves}
	\pause
	\put(22, 51)	{\colorbox{white}{\bl{Pressure}}}
	\put(67, 51)	{\colorbox{white}{\bl{Shear}}}
	\end{overpic}}
\end{frame}

\begin{frame}\frametitle{Vertical \href{https://en.wikipedia.org/wiki/Seismometer}{Seismometer}}
\center
	%\href{https://en.wikipedia.org/wiki/Seismometer\#/media/File\:Seismographs.jpg}{
	\href{https://en.wikipedia.org/wiki/Seismometer}{
	\begin{overpic}[ scale = 0.5 ]
		{\pLocalGraphics seismograph/wiki-01}
		%\put(50, 50)	{\colorbox{white}{$a+b$}}
	\end{overpic}}
\end{frame}

\begin{frame}\frametitle{\href{https://www.iris.edu/hq/inclass/animation/seismograph_vertical_slow_motion}{Seismograph: Vertical (slow motion)}}
\center
	\href{https://www.iris.edu/hq/inclass/animation/seismograph_vertical_slow_motion}{
	\begin{overpic}[ scale = 0.35 ]
		{\pLocalGraphics movie}
		%\put(50, 50)	{\colorbox{white}{$a+b$}}
	\end{overpic}}
\end{frame}

%
%\begin{frame}\frametitle{Test}
%	\movie[ ]{test}{./local/articles/A_004A_seismographverticalslow.mp4}
%\end{frame}
%

%\begin{frame}\frametitle{Vertical \href{https://en.wikipedia.org/wiki/Seismometer}{Seismometer}}
%\center
%	\href{https://www.usgs.gov/media/images/cartoon-sketch-seismograph}{
%	\begin{overpic}[ scale = 0.75 ]
%		{\pLocalGraphics seismograph/seismograph-03}
%		%\put(50, 50)	{\colorbox{white}{$a+b$}}
%	\end{overpic}}
%\end{frame}
%
%\begin{frame}\frametitle{\href{https://www.usgs.gov/faqs/seismometers-seismographs-seismograms-whats-difference-how-do-they-work}{Sample Seismograph}}
%\center
%	\href{https://en.wikipedia.org/wiki/Seismometer\#/media/File\:Digital\_seismogram.png}{
%	\begin{overpic}[ scale = 0.85 ]
%		{\pLocalGraphics seismograph/Digital_seismogram}
%		%\put(50, 50)	{\colorbox{white}{$a+b$}}
%	\end{overpic}}
%\end{frame}

\begin{frame}\frametitle{\href{https://www.usgs.gov/faqs/seismometers-seismographs-seismograms-whats-difference-how-do-they-work}{Sample Seismograph}}
\center
	\href{https://en.wikipedia.org/wiki/Seismometer\#/media/File\:Digital\_seismogram.png}{
	\begin{overpic}[ scale = 0.85 ]
		{\pLocalGraphics seismograph/Digital_seismogram}
		%\put(50, 50)	{\colorbox{white}{$a+b$}}
	\end{overpic}}
	\pause
	\ \\[5pt]
	\href{https://www.dailynews.com/2023/05/08/magnitude-3-4-earthquake-hits-near-coast-in-malibu-area/}{\textrm{The record from station HFP (Fremont Peak) for a 2023 earthquake which occurred 3.7 miles below the surface off the coast of Malibu.}}
\end{frame}

\begin{frame}\frametitle{\href{https://www.iris.edu/hq/inclass/animation/seismograph_vertical_slow_motion}{Waves: S and P}}
\center
	\href{https://en.wikipedia.org/wiki/Seismology}{
	\begin{overpic}[ scale = 0.15 ]
		{\pLocalGraphics seismograph/Seismic_wave_travel_through_Earth}
		%\put(50, 50)	{\colorbox{white}{$a+b$}}
	\end{overpic}}
\end{frame}

\begin{frame}\frametitle{Seismic Wave Paths}
\center
	\href{https://en.wikipedia.org/wiki/File:Earthquake\_wave\_paths.svg}{
	\begin{overpic}[ scale = 0.15 ]
		{\pLocalGraphics seismograph/Earthquake_wave_paths}
		%\put(50, 50)	{\colorbox{white}{$a+b$}}
	\end{overpic}}
\end{frame}

%%   %%   %%   %%   %%   %%   %%   %%   %%
\subsection{Seismic Event Location}
\begin{frame}\frametitle{Hypocenter Location\jumpLittle}
\center
	\href{https://en.wikipedia.org/wiki/Seismic_wave}{
	\begin{overpic}[ scale = 0.35 ]
		{\pLocalGraphics seismograph/Hypocenter_Calculation}
		%\put(50, 50)	{\colorbox{white}{$a+b$}}
	\end{overpic}}
\end{frame}

%%   %%   %%   %%   %%   %%   %%   %%   %%
\subsection{Analysis of Seismic Data}
%
\begin{frame}\frametitle{Fundamentals of Geophysical Data Processing}
\center
	\href{run:../local/articles/Claerbout-1985-Fundamentalsofgeophysicaldataprocesing.pdf}{
	\begin{overpic}[ scale = 0.4 ]
		{\pLocalGraphics claerbout}
		\put(100, 10)	{\href{https://library.seg.org/doi/10.1190/1.1442046}{\scriptsize{\colorbox{gray!10}{\cite{claerbout1986fundamentals}}}}}
	\end{overpic}}
\end{frame}
%
\begin{frame}\frametitle{\href{https://www.cambridge.org/core/books/time-series-analysis-and-inverse-theory-for-geophysicists/AB0EC7F97FFAC6EECA30426879C1987F}{Time Series Analysis and Inverse Theory for Geophysicists}}
\center
	\href{https://www.amazon.com/Time-Analysis-Inverse-Theory-Geophysicists/dp/0521525691}{
	\begin{overpic}[ scale = 0.125 ]
		{\pLocalGraphics gubbins}
		\put(90, 10)	{\href{https://www.cambridge.org/core/books/time-series-analysis-and-inverse-theory-for-geophysicists/AB0EC7F97FFAC6EECA30426879C1987F}{\scriptsize{\colorbox{gray!10}{\cite{gubbins2004time}}}}}
	\end{overpic}}
\end{frame}
%
\begin{frame}\frametitle{\href{https://www.cambridge.org/core/books/time-series-analysis-and-inverse-theory-for-geophysicists/AB0EC7F97FFAC6EECA30426879C1987F}{Time Series Analysis and Inverse Theory for Geophysicists}}
\center
	\href{rhttps://www.amazon.com/Time-Analysis-Inverse-Theory-Geophysicists/dp/0521525691}{
	\begin{overpic}[ scale = 0.55 ]
		{\pLocalGraphics gubbins-02}
		\put(90, 10)	{\href{https://www.cambridge.org/core/books/time-series-analysis-and-inverse-theory-for-geophysicists/AB0EC7F97FFAC6EECA30426879C1987F}{\scriptsize{\colorbox{gray!10}{\cite{gubbins2004time}}}}}
	\end{overpic}}
\end{frame}
%
\begin{frame}\frametitle{Inversion of Seismic Waveforms}
\center
	\href{run:../local/articles/Tarantola-1988-TheoreticalBackgroundfortheInversionofSeismic.pdf}{
	\begin{overpic}[ scale = 0.375 ]
		{\pLocalGraphics tarantola}
		\put(90, -10)	{\href{https://library.seg.org/doi/10.1190/1.1442046}{\scriptsize{\colorbox{gray!10}{\cite{tarantola1986strategy}}}}}
	\end{overpic}}
\end{frame}

\begin{frame}\frametitle{\brune:\spectral}
\center
	\href{run:../local/articles/Brune-1970-Tectonicstressandthespectraofseismicshear.pdf}{
	\begin{overpic}[ scale = 0.5 ]
		{\pLocalGraphics brune-01}
		\put(80, -5)	{\href{\urlbrune}{\scriptsize{\colorbox{gray!10}{\cite{brune1970tectonic}}}}}
	\end{overpic}}
\end{frame}

\begin{frame}\frametitle{\brune:\spectral}
\center
%Spectrum: \\[10pt]
%	\href{run:../local/articles/Brune-1970-Tectonicstressandthespectraofseismicshear.pdf}{
%	\begin{overpic}[ scale = 0.75 ]
%		{\pLocalGraphics brune-03}
%		%\put(80, -5)	{\href{https://agupubs.onlinelibrary.wiley.com/doi/abs/10.1029/JB075i026p04997}{\scriptsize{\colorbox{gray!10}{\cite{brune1970tectonic}}}}}
%	\end{overpic}}
%RMS Average Far-Field Spectrum: \\[10pt]
%	\href{run:../local/articles/Brune-1970-Tectonicstressandthespectraofseismicshear.pdf}{
%	\begin{overpic}[ scale = 0.75 ]
%		{\pLocalGraphics brune-04}
%		%\put(80, -5)	{\href{https://agupubs.onlinelibrary.wiley.com/doi/abs/10.1029/JB075i026p04997}{\scriptsize{\colorbox{gray!10}{\cite{brune1970tectonic}}}}}
%	\end{overpic}}
Far-field spectrum: 
\begin{equation}
	\rd{\Omega}\paren{\bl{\omega}} = \frac{ \sigma \beta }{\mu} \frac{1}{\bl{\omega}\sqrt{\bl{\omega}^{2} + \tau^{-2}}}
\label{eq:average-far-field}
\end{equation}
\ \\[20pt]
RMS average far-field spectrum: 
\begin{equation}
	\left< \rd{\Omega}(\bl{\omega}) \right> = \left<\mathcal{R}_{\theta \phi}\right> \frac{ \sigma \beta }{\mu} \frac{r}{R} F(\epsilon) \frac{1}{\bl{\omega}^{2} + \alpha^{2}}
\label{eq:average-far-field}
\end{equation}
\end{frame}

\begin{frame}\frametitle{\href{https://pubs.geoscienceworld.org/ssa/bssa/article-abstract/63/3/1133/117202/The-spectral-theory-of-seismic-sources}{\randall: Spectral Theory of Seismic Sources}}
\center
	\href{run:../local/articles/randall-spectral-theory.pdf}{
	\begin{overpic}[ scale = 0.45 ]
		{\pLocalGraphics randall-01}
		\put(85, -10)	{\href{\urlrandall}{\scriptsize{\colorbox{gray!10}{\cite{randall1973spectral}}}}}
	\end{overpic}}
\end{frame}

\begin{frame}\frametitle{\href{https://pubs.geoscienceworld.org/ssa/bssa/article-abstract/63/3/1133/117202/The-spectral-theory-of-seismic-sources}{\randall: Spectral Theory of Seismic Sources}}
\center
\bl{Energy} and Magnitude \bl{$M_{\perp}$}\\[10pt]
	\href{run:../local/articles/randall-spectral-theory.pdf}{
	\begin{overpic}[ scale = 0.5 ]
		{\pLocalGraphics randall-02}
		%\put(85, -10)	{\href{https://pubs.geoscienceworld.org/ssa/bssa/article-abstract/70/1/1/101951/A-spectral-theory-for-circular-seismic-sources}{\scriptsize{\colorbox{gray!10}{\cite{randall1973spectral}}}}}
	\end{overpic}}
\end{frame}

\begin{frame}\frametitle{\href{https://en.wikipedia.org/wiki/Spectral_theory}{Boatwright: Spectral Theory}}
\center
	\href{run:../local/articles/bssa0700010001.pdf}{
	\begin{overpic}[ scale = 0.25 ]
		{\pLocalGraphics boatwright}
		\put(80, -5)	{\href{https://pubs.geoscienceworld.org/ssa/bssa/article-abstract/70/1/1/101951/A-spectral-theory-for-circular-seismic-sources}{\scriptsize{\colorbox{gray!10}{\cite{boatwright1980spectral}}}}}
	\end{overpic}}
\end{frame}

\begin{frame}\frametitle{\href{https://en.wikipedia.org/wiki/Spectral_theory}{Boatwright: Spectral Theory}}
\center
	\href{run:../local/articles/bssa0700010001.pdf}{
	\begin{overpic}[ scale = 0.35 ]
		{\pLocalGraphics boatwright-01}
		%\put(-80, 150)	{\href{https://pubs.geoscienceworld.org/ssa/bssa/article-abstract/70/1/1/101951/A-spectral-theory-for-circular-seismic-sources}{\scriptsize{\colorbox{gray!10}{\cite{boatwright1980spectral}}}}}
	\end{overpic}}
\end{frame}


\endinput  %  ==  ==  ==  ==  ==  ==  ==  ==  ==
